\begin{abstract}
Disaggregated LLM serving systems such as \mc pool the DRAM and SSD of
entire GPU clusters into a deduplicated KVCache that every request can reuse,
cutting prefill computation substantially. We show that this sharing is
observable in ways existing defenses do not address. Beyond the known
time-to-first-token channel, a cluster-scale KVCache leaks through cache-aware
admission control, through billed cached-token counts, through the latency
ladder of its storage tiers, and through eviction pressure that survives
per-tenant key salting. Current defenses close such channels only by
partitioning or selectively isolating the cache, forfeiting the reuse that
justifies the architecture.
Our key insight is that confidentiality requires isolating what a principal can
\emph{observe}, not what the system \emph{stores}. \sys keeps the physical
KVCache global and deduplicated, but derives every client-visible
quantity---first-token release time, admission decisions, and usage
accounting---from a per-principal virtual cache with a fixed quota. A
slack-aware scheduler converts the latency withheld from cross-tenant hits into
capacity for other requests. We implement \sys on vLLM and \mc Store and
evaluate it against six attacks using \mc production traces. \sys reduces
all \tbd{six} channels to chance-level distinguishability while retaining
\tbd{$\geq$90\%} of global sharing's prefill savings and serving \tbd{X\%} more
SLO-compliant requests than per-tenant isolation.
\end{abstract}
